% conclusion.tex
\section{Conclusión}\label{sec:conclusion}

\subsection{Conclusiones Principales}

Este estudio implementó y comparó tres métodos de interpolación para el patrón Bayer RGGB:

1. **Interpolación lineal:** El método más simple y rápido, adecuado para sistemas con recursos limitados.
2. **Interpolación bicuadrada:** Ofrece un buen equilibrio entre calidad y complejidad computacional.
3. **Interpolación bicúbica:** Proporciona la mejor calidad visual a costa de mayor complejidad computacional.

Los resultados muestran que cada método tiene su ámbito de aplicación ideal, dependiendo de los requisitos de calidad y las restricciones computacionales del sistema.

\subsection{Contribuciones del Trabajo}

\begin{enumerate}
    \item \textbf{Implementación educativa:} Código MATLAB documentado y modular
    \item \textbf{Análisis comparativo:} Evaluación sistemática de métodos polinomiales
    \item \textbf{Visualización didáctica:} Diagramas explicativos del proceso matemático
    \item \textbf{Documentación completa:} Explicación paso a paso de algoritmos
\end{enumerate}

\subsection{Recomendaciones de Uso}

\begin{table}[h]
\centering
\caption{Guía de selección de método de interpolación}
\begin{tabular}{p{4cm}p{6cm}}
\hline
\textbf{Escenario de aplicación} & \textbf{Método recomendado} \\
\hline
Sistemas embebidos en tiempo real & Lineal \\
Fotografía de consumo & Bicuadrada \\
Video de alta calidad & Bicuadrada adaptativa \\
Fotografía profesional & Bicúbica \\
Procesamiento médico/científico & Bicúbica con post-procesamiento \\
Prototipado rápido & Lineal o Bicuadrada \\
\hline
\end{tabular}
\end{table}

\subsection{Trabajo Futuro}

\begin{itemize}
    \item \textbf{Interpolación adaptativa:} Métodos que ajustan parámetros según contenido local
    \item \textbf{Aprendizaje automático:} Redes neuronales para demosaicing
    \item \textbf{Optimización hardware:} Implementación en FPGA/GPU
    \item \textbf{Patrones alternativos:} Extensiones a X-Trans y otros patrones
    \item \textbf{Corrección de artefactos:} Post-procesamiento para reducir aliasing
    \item \textbf{Benchmarking extensivo:} Comparación con métodos comerciales
\end{itemize}

\subsection{Consideraciones Finales}

La selección del método de interpolación debe basarse en un análisis cuidadoso de los requisitos 
de la aplicación específica, considerando:

\begin{itemize}
    \item Restricciones computacionales y de tiempo
    \item Calidad visual requerida
    \item Características del contenido de la imagen
    \item Recursos hardware disponibles
    \item Costo de implementación y mantenimiento
\end{itemize}

Los métodos presentados en este trabajo proporcionan una base sólida para la implementación 
de sistemas de procesamiento de imágenes que requieren interpolación de patrón Bayer, con 
diferentes niveles de compromiso entre calidad y complejidad.